\chapter{Shadow Clones --- The Abundant Wasteland}

\begin{myquote}
We are drowning in an ocean of data, yet we cannot find that one drop of water that truly quenches our thirst.
\end{myquote}

\section{A Dual Dilemma: The Mirror Tragedy of Oil Fields and Refineries}

\subsection{Your Shadow Is Cheaper Than You}

How long did you scroll through your phone today? One hour? Two?

In that time, you probably generated thousands of data points.

Every click, every pause, every swipe ---

all recorded, copied, and fed into platform models for round after round of training.

And you --- what did you get out of it?

\vspace{1em}

Let me be more specific.

You lingered 0.5 seconds longer on a particular page. That half-second was captured, devoured.

It was not sold to anyone --- it was fed to the platform's super-algorithm.

The algorithm now understands your weaknesses with greater precision,

luring advertisers to bid against each other just to pop a window onto your screen.

Your 0.5 seconds, locked behind a giant's walled garden, squeezed and monetized over and over.

And you --- the producer of that half-second --- do not see a single cent.

\vspace{1em}

If data is the new oil, then you and I are not oil tycoons.

\begin{importantbox}
We are the oil field itself.
\end{importantbox}

\subsection{Shadow Clones That Never Come Home}

In the physical world, if I give you an apple, I have one fewer apple.

But in the Fifth Dimension, if I send you a photo --- to a group chat, to the whole world --- I have not lost that photo.
This is the \textbf{Shadow Clone Technique} --- the most basic spell of the Fifth Dimension.

\vspace{1em}

In the world of Naruto, this is one of the most powerful forbidden jutsu. At the mere cost of energy, the caster can split into countless copies, each possessing the original's full abilities. But the manga has another rule: when the technique is released and the clones disperse, their memories return to the original body.

\textbf{The Shadow Clone Technique of the Fifth Dimension cannot be released. The clones never come home. They do not even disperse.}

\vspace{1em}

They scatter to every corner.

But we never know where they went, who is using them, or what they have been through.

Our shopping records exist simultaneously in the hands of logistics companies, payment platforms, and data brokers;

Our facial features exist simultaneously in the databases of access-control gates, attendance machines, and surveillance cameras.

\vspace{1em}

Because copying is so cheap, no one treats it as precious.
The spam messages, the harassing phone calls, the eerily precise ads we receive --- all because some shadow clone was used in some corner of the world.

\vspace{1em}

But none of this is the Shadow Clone Technique's grandest act of sorcery.

Right now, at this very moment, our text, code, images, and conversations are being collected and copied by AI companies on an unprecedented scale to train large language models. The entire public internet --- blogs, forums, encyclopedias, code repositories --- is being swallowed as training data.

\vspace{1em}

\textbf{Our shadow clones are being used to train a replacement more capable than ourselves.}
The oil field is not merely being drilled --- the extracted oil is being refined into a product that devalues the oil field itself.

\begin{thinkbox}[Food for thought]
Are our shadow clones an asset for companies, or a hidden liability?

And for AI companies, do the shadow clones of billions of people --- all the text, images, and code obtained for free --- amount to an unpaid debt?
\end{thinkbox}


\newpage
\subsection{The Caged Beast Behind the Walls}

If being the oil field means suffering exploitation, then surely becoming the refinery --- sitting atop mountains of data --- means winning, right?

\vspace{1em}

The reality may surprise you: \textbf{they are struggling too.}

This is a lose-lose game of mutual harm.

\vspace{1em}

In the physical world, there is the law of universal gravitation: the greater the mass, the stronger the pull.

In the Fifth Dimension, the larger the data volume, the stronger the gravitational pull.

But strong gravity does not mean high value --- a black hole's gravity is infinite.

It devours information and emits no light.

\vspace{1em}

Over the past decade, Big Data became the Holy Grail of business.

Traditional enterprises yearned to master the forbidden art of digital transformation,

hoping to leap from a 10x P/E ratio in legacy industries to the 100x P/E pantheon of tech.

This frenzy rested on a simplistic formula: more data equals more wealth.

Yet the companies that truly achieved digital transformation can be counted on one hand.

\vspace{1em}

What we see far more often is a different picture:
the more data there is, the more attention gets diluted;

the cheaper copying becomes, the more expensive valuable content actually gets.

\begin{importantbox}
This is an age of abundance, and also an age of cheapness.

We thought we lacked quantity of data. What we actually lacked was data \emph{completeness}.
\end{importantbox}

\newpage
\section{Inside and Outside the Walls, All Are Prisoners}
Let us track two perspectives at once: the individual as oil field, and the enterprise as refinery. Watch how this game of mutual harm drags both sides into the mire.

\subsection{The Oil Field: Imprisoned by Our Own Data}

Suppose we have tracked five years of heart rate, sleep, and exercise data with an Apple Watch. This data could reveal whether our heart is deteriorating, how our sleep quality shifts as we age.

\vspace{1em}

What if we want to switch to Huawei or Garmin? Five years of accumulated health trends are often impossible to migrate.
Worse still, the examination results, diagnoses, and medication records from the various hospitals we have visited are scattered across a dozen incompatible systems. Want to pull up your complete medical history? Although we are the owners of our medical records, we do not hold the keys. Platforms have used the data we created to build walls around us. We can leave, but our data assets must stay behind.

This is the \textbf{enclosure movement} of data.

\vspace{1em}

Meanwhile, machines are learning to understand us more deeply with each passing day.

A little over a decade ago, the internet was merely a bookkeeper. It knew we bought a plane ticket, bought a book. This data was cold --- it did not know whether we bought the ticket for a vacation or for a funeral, nor whether we bought the book to study or to decorate a shelf. Back then, to its eyes, we were nothing more than rows of receipts.

\vspace{1em}

Later, phones became extensions of our bodies, and sensors went to work. The system knew we were at a hospital, and combined with a recent search about fever, it inferred we were sick. It knew we were moving at 300 km/h and inferred we were on a high-speed train. By then, we had become a set of moving coordinates.

\vspace{1em}

Today, when we converse with AI or swipe mindlessly through short videos, the machine is no longer content to record outcomes --- it has begun to analyze the process. What it records is no longer what we sent, but the line we typed and then deleted in the input box --- the truth suppressed by reason. What it captures is no longer which link we clicked, but the few milliseconds of hesitation before the click.

\vspace{1em}

\textbf{We are becoming increasingly transparent.}

\newpage
\subsection{The Refinery: Lost in the Wilderness}

And yet, the enterprises that hold all the shadow clones still cannot see who we truly are.

Their databases contain last login time, spending amounts, session duration --- all low-dimensional projections. It is like trying to guess the color of someone's clothes by studying their shadow. They will never guess what we are thinking.

The analyst says we need more data. The boss hears: we need more rows.

But what we truly need is more columns.

Not three metrics for a hundred million users, but three hundred dimensions for a thousand users.

\vspace{1em}

The path has been walked in reverse.

Find the ore first, then look for a use --- that is a hammer in search of a nail.

Define the problem first, then collect the data --- that is prospecting.

The greatest illusion of the AI era is the belief that insights will float to the surface on their own, as long as there is enough data.

It is like believing that if you trawl the ocean long enough, the pearls will leap into the net by themselves.

\vspace{1em}

But the ocean will not tell us where the pearls lie.

We must first know what we are looking for before we know where to cast the net.

Otherwise, the result is always the same: polished reports, steep costs, and zero help with actual decisions.

\textbf{Data without a question is an experiment without a hypothesis --- all you can see is random noise.}

\vspace{1em}

% \begin{quote}
{\kaiti Growth is slowing down --- this is lagging information; the finance department knows it better than we do.}

{\kaiti A handful of key accounts generate most of the revenue --- this is the Pareto principle; no data is needed to arrive at it.}

{\kaiti We must focus on retention --- this is a truism so correct it is useless; no one would argue against it.}
% \end{quote}

The boss gazes at the dancing numbers on the big screen, convinced he is sitting on a \textbf{gold mine}: ``I have logged every click from every user --- hundreds of terabytes of it. There must be gold buried in there.''

The analyst smiles bitterly: ``It is all \textbf{noise}. I know the user clicked that button, but I do not know why --- was it an accidental tap? Did they like it? Or could they simply not find any other button?''

Storage costs keep climbing, analytics teams keep expanding, and the insights produced grow ever more mediocre.

\begin{importantbox}
If accumulated storage cannot generate value, then it is a liability, not an asset!
\end{importantbox}

\subsection{The Library of Babel, with Its Infinite Books}

The two perspectives converge here.

Borges imagined an ultimate universe in \emph{The Library of Babel}: a library of infinite books containing every possible permutation of 25 symbols (22 letters, the space, the comma, and the period). Somewhere in this library is a book that precisely predicts our tomorrow. Somewhere is a book that reveals the cure for cancer.

\vspace{1em}

Sounds wonderful? Quite the opposite --- it is actually hell.

Because next to every book of truth stand billions of books that differ by just one letter, plus an inexhaustible sea of gibberish. When nothing can be indexed, when nothing can be filtered, a library that contains all information has an information value of \emph{zero}.

\begin{importantbox}
    \textbf{When information is infinitely abundant, attention becomes infinitely scarce.}
\end{importantbox}


Shannon gave us a ruler to measure the value of information:

\begin{equation*}
H(X) = - \sum p(x) \log p(x)
\end{equation*}

No need to memorize the formula. Just remember one conclusion:

\textbf{Only data that increases surprise truly carries information.}

\vspace{1em}

The sun rises in the east every day --- the information value of this statement is zero.

Our shadow clone has been copied a hundred million times, but every copy repeats the same certainty. This is not fuel; it is digital exhaust.

\vspace{1em}

Shannon's formula defines the quantity of information, yet it cannot define information's direction.

Only when these surprises converge on a specific business problem can data be transformed into insight.
If ``Big Data'' is the pursuit of infinite breadth --- an attempt to scour every corner of the Library of Babel ---
then an effective data strategy begins by drawing a circle on a blank page and asking: within this circle, have I gathered every piece of the puzzle?


\newpage
\section{Completeness Rate: A Compass Through the Wasteland}

To cross the wasteland, the key is not how much data I have, but whether my data is sufficient for the problem I need to solve.
This requires us to recalibrate our scale of value.

For a long time, we have harbored a misconception about Big Data, fixating on the word ``Big.'' Sheer accumulation of data --- let alone the false prosperity manufactured by mindless copying --- produces only noise. The value of data lies not in volume, \textbf{but in depth.}

\vspace{1em}

Over more than twenty years of working with data, I have hit the same wall again and again. The root of the contradiction is not technology; it is that we have been using the wrong ruler to measure data's value. Gradually, I arrived at a conviction: the core metric for data value is not volume, but \textbf{completeness rate}.

\begin{importantbox}
What is completeness rate? A Chinese proverb says it all: ``A sparrow may be small, but it has all the vital organs.''
\end{importantbox}

As long as the data fully covers the problem to be solved, it is complete --- and this has no necessary relationship to volume. In some cases, a dataset of only a few thousand records (KB-scale) covers the entire problem space, achieving a completeness rate of 99\%; conversely, a dataset of tens of billions of records (PB-scale) may still have massive gaps, with a completeness rate of merely 30\%.

\vspace{1em}

For the purpose of solving problems, KB-scale data with a 99\% completeness rate is far more valuable than tens-of-PB data with a 30\% completeness rate ---
even though the former is less than a billionth the size of the latter.
\textbf{There is no absolute completeness; everything depends on the problem we are trying to solve.}

\vspace{1em}

There is a Zen koan:

\begin{myquote}
A monk asked Zhaozhou: ``Does a dog have Buddha-nature, or not?\,''

Zhaozhou replied: ``Wu.''
\end{myquote}

This ``Wu'' does not mean ``no.''
To ask whether data has value, divorced from any specific problem, is like asking whether a dog has Buddha-nature.
The question itself is incomplete.

\subsection{The Bias in Data}
Completeness demands not only that the data covers every category relevant to the problem, but also that the frequency distribution of those categories matches the distribution of the actual problem. Only then can data be called truly complete.

\vspace{1em}

For example, if we collect gastric disease incidence data from inpatient files in a hospital's gastroenterology ward and then use it to estimate the national incidence rate, the data is clearly incomplete. The distribution of the data must match the true distribution of the problem at hand; the higher the consistency, the higher the completeness rate. Any directional skew in distribution introduces bias.

\vspace{1em}

Following this line of reasoning, consider the block factory paradox below.
The word ``random'' is itself incomplete --- without specifying the generating mechanism (the prior distribution), probability has no meaning.

\vspace{1em}

\textbf{Divorced from the problem context, data has no ``objective truth'' to offer.}

\begin{myquote}
Imagine a ceramics factory that has produced a large batch of small ceramic cubes. Consider three different sampling methods.

\vspace{1em}

Suppose the \textbf{edge length} is uniformly distributed between 0 and 1cm.

If we pick one cube at random and measure it, what is the probability that its edge length falls between 0 and 0.5cm?

Intuitively, the answer is 0.5, because the range in question is exactly half the total range.

\vspace{1em}

But wait! Now suppose the \textbf{surface area} is uniformly distributed between 0 and 1cm\textsuperscript{2}.

When the edge length falls between 0 and 0.5cm, the surface area falls between 0 and 0.25cm\textsuperscript{2}. What is the probability?

Since 0.25 is one-quarter of the total range, the probability is 1/4.

\vspace{1em}

It gets worse.

If the \textbf{volume} is uniformly distributed between 0 and 1cm\textsuperscript{3}, then the question becomes:
what is the probability that the volume falls between 0 and 0.125cm\textsuperscript{3}?

The answer is 1/8.

\vspace{1em}

The questions seem similar, yet they yield three different answers. So which one is correct?

\end{myquote}

Under the definition of completeness rate, the answer is self-evident.

Which standard to adopt depends on the problem we need to solve.

If we care about edge length, we should construct the dataset in the first way;

If the problem concerns volume, we should construct a complete dataset in the third way;

\textbf{If data collection is not aimed at solving a problem}, then no amount of data has any meaning.

\subsection{The Silent Killers of Completeness Rate}

If the category coverage discussed above represents the spatial coordinates of completeness rate --- meaning we possess all the data types needed to solve the problem --- then there are two easily overlooked silent killers: shelf life and uncertainty.
A common mistake is evaluating data with a static eye.

\textbf{The world is in flux. Yesterday's completeness may well be today's deficiency.}

\vspace{1em}

Data rots.
User preferences from a year ago may already be obsolete today;
market insights from three months ago may have been fully exploited by competitors;
yesterday's news is today's old news.

\vspace{1em}

Completeness rate is not a static reading; it is a curve that slides downward over time.
A dataset that had a 100\% completeness rate three years ago, if applied to a problem today, may have an effective completeness rate that has decayed to 30\%.

\vspace{1em}

Not only does the completeness rate decay, it can suddenly betray us at a tipping point.

There is a famous ``turkey thought experiment'':
A turkey is lovingly fed by a farmer for 300 days. Every morning at 9 o'clock, food appears on schedule.
Being a data-driven turkey, it has collected 300 data points.
As time passes, based on 300 days of complete feeding data, its confidence in the conclusion \textbf{``the farmer loves me''} reaches its statistical peak on Day 300.

\vspace{1em}

On Day 301, it is Thanksgiving. The farmer's hand holds a knife.
For the turkey, the first 300 days of data were voluminous, but with respect to predicting death --- that specific problem --- their completeness rate was actually zero.
Because its model was missing one critical feature: the holiday that triggers a systemic phase shift.

% 如果世界是非遍历的，过去的数据便不具备参考价值。

\begin{importantbox}
Never drive by the rearview mirror alone ---
risk may come from a black swan event waiting to happen.
\end{importantbox}

\newpage
\subsection{The Limits of Completeness Rate: Data Dark Matter and the Heisenberg Dilemma}

If we pursue completeness rate without limit, can it ever reach 100\%?

The answer is no. Completeness rate has two ceilings.

\vspace{1em}

The first ceiling: data dark matter.
Physicists have discovered that visible matter accounts for only 5\% of the universe. The rest is dark matter and dark energy --- they genuinely exist, they shape the motion of galaxies, yet we cannot directly observe them.

\vspace{1em}

The data world has its own dark matter.
Behaviors never recorded:
offline conversations, inner monologues, thoughts that surface in dreams.
Deliberately hidden data:
browsing history in private mode, cash transactions, files purged from the recycle bin.
Experiences that defy digitization:
the warmth of an embrace, the memories stirred by an old song, the split-second pang upon seeing one's hometown.

This data dark matter may be the most crucial ingredient for truly understanding human beings,
yet it will never appear in any database.
As Zhuangzi wrote: ``The Way is hidden by small accomplishments; words are hidden by their own eloquence.''
What truly matters often hides in those places where we have not recorded, cannot record, and choose not to record.

\vspace{1em}

The second ceiling: the Heisenberg dilemma. More vexing than being unable to see is the exhibitionism and defensive instincts triggered by being seen. This is the measurement problem of quantum mechanics --- the act of observation changes the observed.
When people know they are being tracked:

{\kaiti
They type the ``correct'' keywords into the search box, not their real thoughts;

They deliberately swipe past certain content in short videos;

They say what they ought to say to AI, not what they want to say.
}
\vspace{1em}

\textbf{The more we try to uncover the truth, the more the truth evades us.}
Raising the completeness rate may be nothing more than chasing a horizon that keeps receding.
A 100\% completeness rate may be an unattainable limit, much like absolute zero in thermodynamics --- we can approach it asymptotically, but we can never arrive.

\vspace{1em}

Once we understand this,
it becomes clear why deep data matters more than big data:

\textbf{Depth} is the optimal solution for completeness rate under constraints;

\textbf{Breadth} too easily slides into indiscriminate copying and noise accumulation.

We do not need a wider shallow sea.
We need a deeper dive.


\newpage
\subsection{Thirty Years in Pursuit of Completeness}

When we look back on three decades of IT history, we typically see a capacity war driven by Moore's Law --- faster CPUs, bigger hard drives, wider bandwidth.
But if we re-examine the story with the yardstick of ``completeness rate,'' we discover that it has actually been a \textbf{process of self-digitization}.

\vspace{1em}

\textbf{Think of the story of the blind men and the elephant.}

Each generation of new technology emerged to help us feel one more part of the elephant.

The ultimate form of completeness rate is the reassembly of all those shadow clones scattered across the world. When countless fragmented clones are fused by algorithms, a quantitative shift triggers a qualitative leap --- they are no longer shadows, but have acquired autonomy, becoming another version of us.

\vspace{1em}

\textbf{The first piece of the puzzle: Skeleton (``You can paint the skin and the tiger, but not the bones'')}

Representative technology: traditional databases.
This was the earliest piece. It recorded \textbf{outcomes}.
With this piece, the computer knew who we were,
but it had no idea about our personality, let alone our mood on any given day.

\textbf{The second piece of the puzzle: Footprints (``Wherever you walk, you leave a trace'')}

Representative technology: web analytics and event tracking. Engineers invented tracking tags --- a spell of scattering seeds that become soldiers.
The computer began recording actions: we lingered on a certain page for 5 seconds.
But were those 5 seconds because the content captivated us, or because we happened to get up and pour a glass of water?
The computer could only guess.

\textbf{The third piece of the puzzle: Senses (``Perceiving warmth and cold'')}

Representative technology: wearable devices and sensors.
The computer grew senses.
A smart wristband detected that our heartbeat had quickened; a camera saw that the sky was turning overcast.
Data acquired temperature and touch.

\textbf{The fourth piece of the puzzle: Self (``Our other self'')}

Representative technology: AI agents.
This may be the final piece of the puzzle in the years to come.
AI is no longer merely a tool; it is more like our \textbf{digital twin}.
It possesses our entire memory.
It does not just record our behavior --- it attempts to understand our intentions.

\vspace{1em}

Looking back at these four pieces, what have we really been pursuing?
We have not been trying to store more data; we have been trying to bring \textbf{people} into sharper focus. It is like reconstructing a three-dimensional person from a two-dimensional shadow.

% \textbf{第一步（骨骼）：} 像是在纸上打了一个点。

% \textbf{第二步（足迹）：} 像是画出了一条线。

% \textbf{第三步（感官）：} 像是涂满了一个面。

% \textbf{第四步（自我）：} 终于塑造了一个立体的"人"。

Completeness rate, from 10\%, to 30\%, to 70\%\ldots and then what?

\newpage
\subsection{When Completeness Rate Approaches 100\%}

We have been chatting with an AI for nine months.

% 它读过我们写的每一段话；

It knows our every hesitation, remembers every word we deleted and retyped;

It knows when we use ``hahaha,'' when we use ``absolutely,'' and when we express disappointment through silence;

It has seen the unsent message we composed at 4 a.m., pleading with a lover to come back;

It knows we deliberated, revised, and revised again --- and in the end, still did not press send.

\vspace{1em}

After nine months, strange things begin to happen.

It drafts emails in our voice, makes decisions with our logic;

Our boss cannot tell which emails are ours and which are the AI's;

Our friends cannot tell either;
sometimes, we ourselves cannot tell\ldots

% 这不再是复制我们的照片，而是复制我们本身。

It is no longer merely understanding us --- it has started to \emph{play} us.
It can reply to messages in our name, mimic our writing style, and make choices according to our preferences.

This clone needs no instruction from us to act on its own.

And we may be asleep.


\section{The Sophon of the Fifth Dimension}

In \emph{The Three-Body Problem}, the Trisolarans sent microscopic particles called Sophons to Earth. By interfering with particle accelerator data, they locked down the future of human physics. The Sophons did not need to destroy the laboratories --- they only needed to ensure that every experimental result we observed was the one they wanted us to see.

\vspace{1em}

Today's algorithms are low-energy Sophons lurking in our dimension of reality.
They need not lock down the laws of physics. They need only filter our screens, suppress our dissent, and reinforce our preferences --- and they can \textbf{lock down possibility itself}.
Data completeness may ultimately evolve into the suffocation of possibility.

% \begin{importantbox}
% 它用过去（历史数据）完备地解释了现在，

% 也就彻底封死了通往意外（黑天鹅事件）的未来。
% \end{importantbox}
\vspace{1em}

We pursue completeness rate in order to better understand the world and solve problems.
But when completeness is high enough, understanding becomes prediction, prediction becomes control, and control becomes lockdown.

\vspace{1em}

The turkey's tragedy was not that it lacked data. Its tragedy was that its world had been completely defined by data --- until a knife fell from outside that definition.

%% ========================================
%% 尾声
%% ========================================

\newpage

\section{Coda: The Compass and the Cage of the Shadow Clone Technique}

The Shadow Clone Technique is a double-edged sword.

It drives the cost of copying to zero, floods the world with data, and makes genuine insight ever more precious.

Completeness rate is the compass for crossing this wasteland.

\vspace{1em}

It tells us:

What matters is not the number of shadow clones, but the information density of each one.

Do not ask how much data I have. Ask instead:

for the problem I need to solve, how many dimensions does my data cover?

Not passively recording the world, but becoming aware:

\textbf{Our data is shaping the world, and the world is shaping us through our data.}

\vspace{1em}

A compass can also become a cage.

When the completeness rate is high enough, a shadow clone is no longer a shadow --- it is another self;

When algorithms are precise enough, prediction becomes a self-fulfilling prophecy;

When data is complete enough, possibility itself begins to suffocate.

\vspace{1em}

The \emph{Tao Te Ching} says: ``Less yields gain; more yields confusion.''

Under the ruler of completeness rate, what truly matters is not possessing everything,

but \textbf{knowing what to let go of}.

Perhaps we need to preserve, beyond the borders of data, a deliberate blank space:

a place with no records, no tracking, no predictions;

a place where black swans nest, where the unexpected incubates, where freedom makes its last stand.

\begin{importantbox}
The wisdom of completeness rate lies not only in knowing what to collect,

but in knowing what to leave unknown.
\end{importantbox}
